\subsection{集合族乘积的定义}

让 \((\mathbf{X}_{\iota})_{\iota \in \mathbf{I}}\) 是一族集合，且设 \(\mathbf{F}\) 是一个定义域为 \(\mathbf{I}\) 使得 \(\mathbf{F}(\iota) \in \mathbf{X}_{\iota}\) 对于每一个 \(\iota \in \mathbf{I}\) 然后对每一个 \(\iota \in \mathbf{I}\) 我们拥有 \(\mathbf{F}(\iota) \in \mathbf{A} = \bigcup_{\iota \in \mathbf{I}} \mathbf{X}_{\iota}\) ，从而 \(\mathbf{F}\) 是 \(\mathfrak{P}(\mathbf{I} \times \mathbf{A})\) 因此，具有上述性质的函数图构成以下集合的一个子集 \(\mathfrak{P}(\mathbf{I} \times \mathbf{A})\) .

定义1. 设 \((\mathbf{X}_{\iota})_{\iota \in \mathbf{I}}\) 是一个集合族。函数图 \(\mathbf{F}\) （以 I 为定义域，使得 \(\mathrm{F}(\iota) \in \mathbf{X}_{\iota}\) 对于每一个 \(\iota \in \mathbf{I}\) ）组成的集合，称为该集合族 \((\mathbf{X}_{\iota})_{\iota \in \mathbf{I}}\) 的乘积，并且记作 \(\prod_{\iota \in \mathbf{I}} \mathbf{X}_{\iota}\) 。对于每一个 \(\iota \in \mathbf{I}\) , \(\mathbf{X}_{\iota}\) 称为指标为 \(\iota\) 的因子（在乘积 \(\prod_{\iota \in \mathbf{I}} \mathbf{X}_{\iota}\) 中）。映射 \(\mathrm{F} \to \mathrm{F}(\iota)\) ( \(\mathrm{F} \in \prod_{\iota \in \mathbf{I}} \mathbf{X}_{\iota}\) , \(\mathrm{F}(\iota) \in \mathbf{X}_{\iota}\) ）被称为指标为 \(\iota\) 的坐标函数（或投影），并记作 \(\operatorname{pr}_{\iota}\) .

\(F(t)\) 被称为 F 的指标为 \(t\) 的坐标（或指标为 \(t\) 的投影）；像 \(pr_{t}\langle A\rangle\) （ \(\prod_{i\in I}X_{i}\) 的子集 A 在指标为 \(t\) 的坐标函数下的像）称为 A 的指标为 \(t\) 的投影。容易验证 \(A\subset\prod_{i\in I}pr_{t}\langle A\rangle\) .

我们常使用记号 \((x_{i})_{i\in I}\) 来表示一个元素 \(\prod_{i\in I}X_{i}\) （§3，第6号）。

如果 \(\mathbf{I} = \varnothing\) ，集合 \(\prod_{i\in \mathbf{I}}\mathbf{X}_i\) 只有一个元素，即空集（§3，第4节，例1）。

如果所有因子 \(X_{t}\) 的乘积 \(\prod_{i\in I}X_{t}\) 都等于同一个集合 E，我们有 \(\prod_{i\in I}X_{t}=E^{I}\) ；这直接从定义得出。

如果 \((\mathbf{X}_t)_{t\in I}\) 是任意一族集合，并且如果 \(\mathbf{E}\) 是这样的一个集合，使得

\[
\bigcup_ {i \in I} X _ {i} \subset E,
\]

那么定义1表明 \(\prod_{i\in I}X_i\subset E^I\) ；因此，在 \(\prod_{i\in I}X_i\) 与从 \(I\) 到 \(E\) 的映射组成的一组集合（即 \(\mathcal{F}(I,E)\) 的一个子集）之间存在一个一一对应关系.

如果 \(\mathbf{I} = \{\alpha\}\) 是由单个元素组成的集合，我们有

\[
\prod_ {i \in I} X _ {i} = X _ {\alpha} ^ {\{\alpha \}};
\]

映射 \(\mathbf{F} \to \mathbf{F}(\alpha)\) 则是一个双射（称为典范双射） \(\prod_{i \in \{x\}} \mathbf{X}_i\) 到上 \(\mathbf{X}_\alpha\) .

设 A、B 为集合，且设 \(\alpha, \beta\) 为不同的对象（存在两个不同的对象，例如 \(\phi\) 并且 \(\{\phi\}\) ）。考虑图 \(\{(\alpha, A), (\beta, B)\}\) ，这显然是泛函的，并且正是该族。 \((X_{t})_{t \in [\alpha, \beta]}\) 使得 \(X_{\alpha} = A\) 并且 \(X_{\beta} = B\) 对于每一对 \((x, y) \in A \times B\) ，让 \(f_{x,y}\) 是函数图 \(\{(\alpha, x), (\beta, y)\}\) 。显然，函数 \((x, y) \to f_{x,y}\) 是 \(A \times B\) 到上 \(\prod_{t \in [\alpha, \beta]} X_t\) ；这个双射的逆映射是 \(g \to (g(\alpha), g(\beta))\) 。这两个双射被称为典范的。在下文中，我们将利用这一一对应关系，从一族集合的乘积的性质推导出两个集合的乘积的性质。

让 \((\mathbf{X}_i)_{i\in I}\) 是一族集合，其中每个集合都只含有一个元素，设为 \(\mathbf{X}_i = \{a_i\}\) ；那么乘积 \(\prod_{i\in I}\mathbf{X}_i\) 是由单一元素组成的集合 \((a_i)_{i\in I}\) .

设 E 是一个集合。将 I 中元素映到 E 中的常值映射 \(\iota \to x\) 的图组成子集 \(\Delta\) （乘积 \(\mathbf{E}^{\mathrm{I}}\) 的一个子集），称为对角线。如果 \(\overline{x}\) 表示映射 \(\iota \to x\) （其中 \(x \in \mathbf{E}\) ）的图，那么映射 \(x \to \overline{x}\) 是从 E 到 \(\Delta\) 上的双射，称为对角映射。

命题4。设 \((\mathbf{X}_i)_{i\in I}\) 是一族集合，并设 \(u\) 是从一个集合 \(K\) 到索引集 \(I\) 上的双射。如果 \(U\) 是 \(u\) 的图，那么映射 \(\mathbf{F} \to \mathbf{F} \circ U\) （从 \(\prod_{i\in I} X_i\) 到 \(\prod_{x\in K} X_{u(x)}\) ）是双射的。

让

\[
\mathrm{A} = \bigcup_ {i \in \mathbf {I}} \mathrm{X} _ {i} = \bigcup_ {\boldsymbol {x} \in \mathbf {K}} \mathrm{X} _ {u (\boldsymbol {x})}
\]

的图像（§4，命题1）。映射 \(\mathbf{F} \to \mathbf{F} \circ \mathbf{U} (\mathbf{F} \in \mathbf{A}^{\mathbf{I}})\) 是 \(\mathbf{A}^{\mathbf{I}}\) 到上 \(\mathbf{A}^{\mathbf{K}}\) （命题2）。显然，“对每个 \(\iota \in \mathbf{I}, \mathbf{F}(\iota) \in \mathbf{X}_{\iota}\) ”的条件等价于“对每个 \(x \in K, (\mathbf{F} \circ \mathbf{U})(x) \in X_{u(x)}\) ”，结论随之成立。

